\section{课程定位与核心命题}

本讲是 Berkeley ``Advanced LLM Agents'' 春季课程中一节高密度专题，由 Google DeepMind 的 Charles Sutton 主讲。讲座看似覆盖两个方向（Code Agent 与 AI Security），但真实主线只有一条：\textbf{如何把语言模型从“会写代码”推进到“能完成开放式软件任务”}。演讲结构按三段展开：先回到 coding agent 的评测与系统设计，再讨论 AI 在安全任务中的适配，最后用 Big Sleep 项目说明在真实漏洞挖掘里 Agent 化方法为什么成立。

\begin{importantbox}{全讲的工程判断标准}
Sutton 的立场并不是 ``模型参数越大越好''，而是 ``系统能否闭环''。闭环包括：可执行任务定义、可验证结果、可迭代轨迹、可扩展工具接口。只要这四项里有一项缺失，Agent 结果就容易停留在演示层面。
\end{importantbox}

演讲前半段多次强调一个容易被忽略的事实：Agent 研发不是单次 prompt 优化，而是\textbf{在巨大组合空间里做 hill climbing}。当团队不断改 system instruction、工具形态、反馈格式、后训练配比时，每一步都依赖评测分数作为导向。若评测目标错位，迭代方向就会整体偏航。

\begin{knowledgebox}{跨领域背景为什么重要}
主讲人先讲自己的经历（NLP $\rightarrow$ 程序合成 $\rightarrow$ 代码生成 $\rightarrow$ 安全），不是铺垫个人履历，而是为了说明：AI 技术价值取决于你对问题域的理解深度。仅懂模型、不了解软件工程或安全流程，最终会把系统优化成 ``会答题''，而不是 ``会做事''。
\end{knowledgebox}

\begin{warningbox}{不要把本讲理解成“更多安全 benchmark”}
讲座的重点不是再介绍几个排行榜，而是解释为什么很多 benchmark 在早期有效、后期失效，以及为什么安全场景要求更强的动态验证。若只记住数据集名字，会错过最有价值的工程方法论。
\end{warningbox}

\begin{figure}[H]
\centering
\includegraphics[width=0.88\textwidth]{frame-01-agenda.jpg}
\caption{讲座早段给出的三段式议程：Code Agent 基础、AI Security、Big Sleep 案例 \protect\footnotemark}
\end{figure}
\footnotetext{视频画面时间区间：00:03:22--00:03:50。}

\subsection{本章小结}
本章建立了读整节课的坐标系：这不是零散技术盘点，而是一堂关于 ``Agent 系统怎样被正确评估、正确约束、正确落地'' 的方法课。后续所有细节都服务于这个核心问题。

